% Generated draft: requires mathematical and source review.
\documentclass[11pt]{article}
\usepackage[margin=25mm]{geometry}
\usepackage{fontspec}
\setmainfont{Latin Modern Roman}
\usepackage{amsmath,amssymb,mathtools,unicode-math}
\setmathfont{Latin Modern Math}
\usepackage{enumitem,needspace,xurl,hyperref}
\setlength{\parindent}{0pt}
\setlength{\parskip}{5pt}
\setlength{\emergencystretch}{3em}
\newcommand{\N}{\mathbb{N}}
\newcommand{\Z}{\mathbb{Z}}
\newcommand{\Q}{\mathbb{Q}}
\newcommand{\R}{\mathbb{R}}
\newcommand{\C}{\mathbb{C}}
\newcommand{\F}{\mathbb{F}}
\newcommand{\A}{\mathbb{A}}
\newcommand{\PP}{\mathbb{P}}
\newcommand{\E}{\mathbb{E}}
\newcommand{\eps}{\varepsilon}
\newcommand{\dd}{\,\mathrm{d}}
\begin{document}
\section*{UnsolvedMath ID: 6200107}
\textbf{Source problem code:} catalogue entry
\subsection{Definitions and mathematical statement}
The supplied question, called the ``Coarse Whitehead Conjecture'', is: under unspecified preceding assumptions, does \(Y\) have coarsely trivial \(\pi_m\) for every integer \(m\ge2\)? Neither \(Y\), its metric or complex structure, nor those assumptions is included in the record. Thus the conditional target can only be retained as
\[
 \text{the missing source hypotheses on }Y
 \quad\Longrightarrow\quad
 \forall m\ge2:\ \pi_m(Y)\text{ is coarsely trivial in the source's sense}.
\]
It must not be read as a statement about every metric space or as ordinary asphericity.

For orientation, one common metric formulation uses the Rips complex \(\operatorname{Rips}_r(Y)\), whose simplices are finite subsets of diameter at most \(r\). A possible uniform triviality condition is
\[
 \forall r>0\ \exists R\ge r\ \forall y\in Y:\quad
 \pi_m(\operatorname{Rips}_r(Y),y)\longrightarrow
 \pi_m(\operatorname{Rips}_R(Y),y)
 \text{ is the zero homomorphism}.
\]
The map is induced by inclusion. This is only an illustrative convention: uniform triviality at a larger scale, eventual null-homotopy for each individual map, and other coarse homotopy theories have different quantifiers. The excerpt does not select among them.

The supplied summary calls the question open but gives no missing assumptions, construction of \(Y\), exact coarse-homotopy definition, or resolution reference. The historical problem label and later open assessment are therefore source information rather than a verified complete conjecture or current literature determination.
\subsection{Short English statement}
Under the source's preceding assumptions, should all higher homotopy groups of \(Y\) vanish in the intended coarse sense? The assumptions, space, and coarse-homotopy convention are missing, so the conjecture cannot yet be stated as a complete theorem.
\subsection{Sources}
\begin{itemize}
\item Local registry: \url{data/unsolvedmath/problems.json}, permanent UnsolvedMath ID \(6200107\), source code catalogue entry; supplied statement, background, and research summary, updated 2026-08-06. Source attribution: ``Coarse Whitehead Conjecture'', Misha Kapovich, 2005.
\item Kapovich, \emph{Problems on Boundaries of Groups and Kleinian Groups} (2005), Problem 107, PDF page 26: \url{https://www.math.ucdavis.edu/~kapovich/EPR/problems.pdf}. Supplied normalized PDF-text extraction accessed 2026-07-29; current status marked as needing review.
\item No separate definition or publication supplying the omitted preceding context is included.
\end{itemize}
\end{document}
